%% ============================================================
%% PAC-Bayesian Generalization Bounds for Weight-Decomposed
%% Low-Rank Adaptation
%% CS 229 Project — Revised & Expanded Draft
%%
%% EDITORIAL NOTES are marked with   %% [EDITORIAL]: ...
%% TODO placeholders are marked with  %% [TODO]: ...
%% ============================================================

\documentclass[11pt]{article}
\usepackage[margin=1in]{geometry}
\usepackage{amsmath, amssymb, amsthm}
\usepackage{mathtools}          % for \coloneqq
%\usepackage{bbm}                % for \mathbbm{1} (install texlive-fonts-extra if needed)
\usepackage{hyperref}
\usepackage{cleveref}
\usepackage{xcolor}
\usepackage{tcolorbox}
\usepackage{booktabs}

%% ------ theorem environments --------------------------------
\newtheorem{theorem}{Theorem}[section]
\newtheorem{lemma}[theorem]{Lemma}
\newtheorem{corollary}[theorem]{Corollary}
\newtheorem{proposition}[theorem]{Proposition}
\newtheorem{definition}[theorem]{Definition}
\newtheorem{remark}[theorem]{Remark}
\newtheorem{assumption}[theorem]{Assumption}

\crefname{lemma}{lemma}{lemmas}
\Crefname{lemma}{Lemma}{Lemmas}

%% ------ editorial / todo box --------------------------------
\newcommand{\editorial}[1]{%
  \begin{tcolorbox}[colback=yellow!10, colframe=orange!60, title=\textbf{Editorial Note}]
    #1
  \end{tcolorbox}}
\newcommand{\todosec}[1]{%
  \begin{tcolorbox}[colback=blue!5, colframe=blue!50, title=\textbf{TODO}]
    #1
  \end{tcolorbox}}

%% ------ notation shorthands ---------------------------------
\newcommand{\KL}{\mathrm{KL}}
\newcommand{\vMF}{\mathrm{vMF}}
\newcommand{\Ad}{A_d}
\newcommand{\Ak}{A_k}
\newcommand{\Rhat}{\hat{R}}
\newcommand{\Rhatl}{\hat{R}_\lambda}
\newcommand{\RQ}{R(Q)}
\newcommand{\RQhat}{\hat{R}(Q)}
\newcommand{\Sd}{\mathbb{S}^{d-1}}
\newcommand{\Sk}{\mathbb{S}^{k-1}}
\newcommand{\N}{\mathcal{N}}
\newcommand{\E}{\mathbb{E}}
\newcommand{\R}{\mathbb{R}}
\newcommand{\phikl}{\phi_\lambda}

\title{\textbf{PAC-Bayesian Generalization Bounds for\\
Weight-Decomposed Low-Rank Adaptation}}
\author{Stephen Zhang}
\date{}

\begin{document}
\maketitle

% ============================================================
\begin{abstract}
Parameter-efficient fine-tuning (PEFT) methods are rising in popularity as foundation models increase in size. One paradigm of PEFT is weight decomposition, in which methods decompose neural network weight updates into direction and magnitude components. DoRA~\cite{liu2024dora} and MAP~\cite{si2025map} are representative methods of the weight decomposition paradigm, yet their generalization behavior remains poorly understood. We use PAC-Bayes to derive generalization bounds that explicitly capture the geometric structure of direction and magnitude updates for weight decomposition PEFT.

Our analysis develops non-asymptotic KL bounds for von Mises–Fisher (vMF) distributions that remain valid outside the high-concentration regime, yielding closed-form complexity terms expressed through angular deviations on the hypersphere. From these we derive a data-dependent Gibbs posterior whose minimizer is a principled training objective, showing how PAC-Bayes gives rise to cosine-based regularization.

We establish a comparison theorem showing that column-wise methods
such as DoRA distribute directional complexity across local angles, while
global methods such as MAP concentrate directional complexity into a single rotation. From the comparison, we theorize DoRA achieves strictly tighter generalization bounds under local sparse updates, whereas MAP is
better under globally aligned updates.

We formalize four testable predictions from the theory and fine-tuning RoBERTa-base across six GLUE tasks under DoRA, MAP, and LoRA to validate them in
future work. We present our findings at the bottom of the paper. 
\end{abstract}

% ============================================================
\section{Introduction}

Large foundation models are often adapted to downstream tasks using
parameter-efficient fine-tuning (PEFT) methods, which modify only a small
subset of parameters while preserving most of the pretrained weights.
Among these approaches, recent methods such as DoRA~\cite{liu2024dora} and
MAP~\cite{si2025map} adopt a \emph{weight-decomposition} perspective,
representing parameters in terms of their magnitudes and directions.
Despite the empirical success of these methods, their generalization properties remain poorly understood.

A key challenge lies in understanding how different parameterizations affect
the complexity of adaptation.
DoRA performs \emph{column-wise} updates, allowing each column of a weight
matrix to rotate independently.
MAP applies a \emph{global} transformation that treats the
entire flattened weight matrix as a single vector and rotates it as a whole.
While these approaches differ geometrically, there is currently no
theoretical framework that explains when or why one should generalize better than
the other.

\begin{sloppypar}
We address this gap through a PAC-Bayesian analysis of weight-decomposed adaptation.
Our approach models directional updates using von Mises-Fisher (vMF) distributions on the hypersphere and derives generalization bounds that depend explicitly on angular deviations between pretrained and fine-tuned parameters.
This leads to a notion of \emph{directional complexity}, measured by cosine-based penalties that arise directly from KL divergence terms.
\end{sloppypar}

%% [EDITORIAL]: Add 1--2 sentences here motivating \emph{why} vMF is the
%% right choice for directional priors.  A useful angle: vMF is the
%% maximum-entropy distribution on $\mathbb{S}^{d-1}$ subject to a fixed
%% mean direction, analogous to the Gaussian being maximum-entropy on
%% $\mathbb{R}^d$ with fixed mean.  This justifies the choice
%% information-theoretically rather than by convenience.

Our analysis yields several key insights.
First, we derive non-asymptotic KL bounds for vMF distributions
(\cref{thm:pac_dora,thm:pac_map}), enabling rigorous treatment of
directional uncertainty without relying on high-concentration
approximations.
Second, we introduce a data-dependent Gibbs posterior
(\cref{sec:gibbs}) that yields a principled training objective,
connecting PAC-Bayes theory to commonly used regularization schemes.
Third, we establish a comparison theorem (\cref{thm:comparison})
showing that column-wise (local) methods distribute angular
changes across many coordinates, while global methods compress them
into a single rotation.

Our results lead to a sharp characterization of when each method is
preferable.
We show that DoRA achieves strictly tighter generalization bounds in regimes
where updates are sparse and localized (\cref{thm:dora_advantage,rem:comparison}),
whereas MAP is favored when updates are globally aligned.
Through our work, we provide the first theoretical explanation for empirical phenomena
observed across PEFT methods.

Finally, we empirically validate our theory by measuring PAC-Bayes
complexity in practical fine-tuning settings (\cref{sec:experiments}).
We show that angular complexity strongly correlates with generalization
performance and that the predicted regimes---sparse versus dense
updates---align with observed performance differences.


%% [TODO]: Add a ``Related Work'' section (see below) before the
%% technical sections; reviewers will expect it.

% ============================================================
\section{Background and Related Work}

\subsection{Parameter-Efficient Fine-Tuning}
Low-Rank Adaptation (LoRA) [3] restricts weight updates to $\Delta W = BA$ via low-rank matrices $B \in \mathbb{R}^{d \times r}$ and $A \in \mathbb{R}^{r \times d}$ ($r \ll d$). Recent approaches enhance this by decoupling magnitude and direction. DoRA [1] factors weight columns as $w_j = m_j u_j$, applying LoRA to the unit-norm direction $u_j \in \mathbb{S}^{d-1}$ while directly training the scalar magnitude $m_j \in \mathbb{R}_+$. Conversely, MAP [2] flattens the entire weight matrix $W \in \mathbb{R}^{d \times d}$ into $w \in \mathbb{R}^{d^2}$, learning global scalars $\alpha, \beta \in \mathbb{R}$ to interpolate between unit-norm vectors via $w' = \alpha \hat{w} + \beta \Delta \hat{w}$.

\subsection{PAC-Bayesian Generalization Theory}
PAC-Bayes bounds [4, 5] provide high-probability generalization guarantees for randomized predictors, driven primarily by the KL divergence $\text{KL}(Q\|P)$ between a learned posterior $Q$ and a data-free prior $P$. While existing literature readily applies PAC-Bayes to standard fine-tuning [6, 7], no prior work analyzes weight-decomposed adaptation.

\subsection{Distributions on Hyperspheres}
The von Mises-Fisher distribution $\text{vMF}(\mu, \kappa)$ on $\mathbb{S}^{d-1}$ is defined by the density:
\[
p(x \mid \mu, \kappa) = C_d(\kappa) \exp(\kappa \mu^\top x), \quad C_d(\kappa) = \frac{\kappa^{d/2-1}}{(2\pi)^{d/2} I_{d/2-1}(\kappa)},
\]
where $I_\nu$ is the modified Bessel function of the first kind. The concentration parameter $\kappa$ controls the distribution, spanning from uniform ($\kappa = 0$) to a point mass at the mean direction $\mu$ ($\kappa \to \infty$). The mean resultant length, $A_d(\kappa) = I_{d/2}(\kappa) / I_{d/2-1}(\kappa) \in [0, 1)$, increases monotonically with $\kappa$.

%% [EDITORIAL]: Originality assessment (see end-of-paper notes).
%% The key gap you are filling is: (1) vMF-based PAC-Bayes with
%% non-asymptotic KL, (2) a geometric comparison between column-wise vs.\
%% global decomposition, and (3) the Gibbs objective interpretation of
%% cosine regularization.  None of these appear in existing literature.

% ============================================================
\section{PAC-Bayesian Analysis of Weight-Decomposed Adaptation}
\label{sec:pac_analysis}

\subsection{Setup and McAllester's Bound}
\label{sec:setup}

Let $\mathcal{H}$ be a hypothesis space and $\mathcal{D}$ be a data
distribution over $\mathcal{Z}$.
Let $P$ be a prior over $\mathcal{H}$ chosen independently of the data,
and $Q$ a posterior learned from a sample $S = \{z_1, \ldots, z_n\} \sim
\mathcal{D}^n$.
Assume the loss $\ell : \mathcal{H} \times \mathcal{Z} \to [0,1]$.
Define the population and empirical risks:
\[
  R(Q) \coloneqq \E_{h \sim Q}\E_{z \sim \mathcal{D}}[\ell(h,z)],
  \qquad
  \hat{R}(Q) \coloneqq \E_{h \sim Q}\frac{1}{n}\sum_{i=1}^n \ell(h, z_i).
\]

\begin{theorem}[McAllester \cite{mcallester1999pac}]
\label{thm:mcallester}
For any $\delta \in (0,1)$, with probability at least $1-\delta$ over the
draw of $S$, for all posteriors $Q$:
\[
  R(Q) \leq \hat{R}(Q) + \sqrt{\frac{\KL(Q \| P) + \ln\frac{2\sqrt{n}}{\delta}}{2n}}.
\]
\end{theorem}
\noindent The problem of bounding $R(Q)$ reduces to bounding $\KL(Q \| P)$
under DoRA and MAP.

% ---------------------------------------------------------------
\subsection{Non-Asymptotic KL Bounds for vMF Distributions}
\label{sec:nonasymptotic}

We drop the high-concentration assumption and work with exact,
non-asymptotic bounds. Lemmas 3.2 and 3.3 are well known.

\begin{lemma}[vMF KL Divergence with Equal Concentration]
\label{lem:vmf_kl}
Let $Q_1 = \vMF(\mu_1, \kappa)$ and $Q_0 = \vMF(\mu_0, \kappa)$ be
distributions on $\Sd$ with the same concentration $\kappa > 0$.
Then:
\[
  \KL(Q_1 \| Q_0) = \kappa \Ad(\kappa)\,(1 - \mu_1^\top \mu_0),
\]
where $\Ad(\kappa) = I_{d/2}(\kappa) / I_{d/2-1}(\kappa)$.
\end{lemma}

\begin{lemma}[Uniform Bounds on $\Ad(\kappa)$]
\label{lem:ad_bounds}
For all $\kappa > 0$ and $d \geq 2$:
\[
  \frac{\kappa}{d + \kappa} \leq \Ad(\kappa) \leq 1.
\]
\end{lemma}

\begin{theorem}[Non-Asymptotic Directional KL Bound]
\label{thm:nonasymptotic}
Under the column-wise decomposition of DoRA with per-column vMF priors
of concentration $\kappa$:
\[
  \frac{\kappa^2}{d+\kappa} \sum_{j=1}^d (1 - \cos\theta_j)
  \;\leq\; C^{\mathrm{DoRA}}_{\mathrm{dir}}
  \;\leq\; \kappa \sum_{j=1}^d (1 - \cos\theta_j).
\]
Under the global decomposition of MAP:
\[
  \frac{\kappa^2}{k+\kappa}\,(1 - \cos\Theta_{\mathrm{global}})
  \;\leq\; C^{\mathrm{MAP}}_{\mathrm{dir}}
  \;\leq\; \kappa\,(1 - \cos\Theta_{\mathrm{global}}).
\]
\end{theorem}

\begin{proof}
Apply \cref{lem:vmf_kl} to get $C^{\mathrm{DoRA}}_{\mathrm{dir}}
= \sum_{j=1}^d \KL(Q_{\mathrm{dir},j} \| P_{\mathrm{dir},j})
= \kappa \Ad(\kappa) \sum_{j=1}^d (1 - \cos\theta_j)$.
The upper and lower bounds follow by substituting the bounds on
$\Ad(\kappa)$ from \cref{lem:ad_bounds}:
\[
  \frac{\kappa^2}{d+\kappa} \sum_{j=1}^d (1-\cos\theta_j)
  \;\leq\;
  \kappa\Ad(\kappa)\sum_{j=1}^d(1-\cos\theta_j)
  \;\leq\;
  \kappa\sum_{j=1}^d(1-\cos\theta_j).
\]
The MAP case is identical, replacing $d$ with $k$ (the ambient dimension
of the flattened weight vector) and the sum with the single global term.
\end{proof}

\begin{remark}
These bounds hold for all $\kappa > 0$, removing the need for
$\kappa \to \infty$ and making the analysis valid in
moderate-concentration regimes.
In particular, when $\kappa$ is small (diffuse prior), the lower and
upper bounds are both $\asymp \kappa^2 / (d-1)$, showing the bounds
are tight up to constants.
\end{remark}

% ---------------------------------------------------------------
\subsection{Column-Wise Decomposition: DoRA}
\label{sec:dora}

Consider a weight matrix $W \in \R^{d \times d}$ decomposed column-wise:
\[
  w_j = m_j u_j, \quad m_j \in \R_+,\; u_j \in \Sd, \quad j = 1,\ldots,d.
\]
We make the following modeling assumptions.

\begin{enumerate}
  \item[(A1)] \textbf{Mean-field factorization.}
    $P = \prod_{j=1}^d P_j$, $Q = \prod_{j=1}^d Q_j$.
  \item[(A2)] \textbf{Independent magnitude and direction.}
    $P_j = P_{\mathrm{mag},j} P_{\mathrm{dir},j}$ and similarly for $Q_j$.
  \item[(A3)] \textbf{Gaussian magnitudes.}
    $P_{\mathrm{mag},j} = \mathcal{N}(m_{0,j}, \sigma^2)$,
    $Q_{\mathrm{mag},j} = \mathcal{N}(\hat{m}_j, \sigma^2)$.
  \item[(A4)] \textbf{vMF directions with shared concentration.}
    $P_{\mathrm{dir},j} = \vMF(\mu_{0,j}, \kappa)$,
    $Q_{\mathrm{dir},j} = \vMF(\hat{\mu}_j, \kappa)$.
\end{enumerate}
%% [EDITORIAL]: Assumption (A3) fixes equal variance between prior and
%% posterior.  This is a simplifying assumption; a reviewer may ask about
%% the case $Q_{\mathrm{mag},j} = \mathcal{N}(\hat{m}_j, \sigma_j^2)$ with
%% learned variance.  You might add a brief remark noting that extending
%% to learned variance only adds a $\frac{1}{2}\ln(\sigma^2/\sigma_j^2) +
%% \frac{\sigma_j^2}{2\sigma^2} - \frac{1}{2}$ term per column, which does
%% not change the qualitative comparison.
Under our assumptions, the KL decomposes as:
\begin{align}
  \KL(Q \| P)
  &= \sum_{j=1}^d \KL(Q_j \| P_j) \notag\\
  &= \sum_{j=1}^d \bigl[\KL(Q_{\mathrm{mag},j} \| P_{\mathrm{mag},j})
     + \KL(Q_{\mathrm{dir},j} \| P_{\mathrm{dir},j})\bigr].
\label{eq:kl_decomp}
\end{align}
The first equality uses (A1) and the product structure. The second equality uses (A2) and the fact that KL factorizes over independent components.
For the Gaussian magnitude terms under (A3):
\begin{equation}
  \KL\bigl(\mathcal{N}(\hat{m}_j, \sigma^2) \;\|\; \mathcal{N}(m_{0,j}, \sigma^2)\bigr)
  = \frac{(\hat{m}_j - m_{0,j})^2}{2\sigma^2}.
\label{eq:kl_gaussian}
\end{equation}
This follows from the standard formula for Gaussians with equal variance.
For the directional terms, Lemma~\ref{lem:vmf_kl} gives:
\begin{equation}
  \KL\bigl(\vMF(\hat{\mu}_j, \kappa) \;\|\; \vMF(\mu_{0,j}, \kappa)\bigr)
  = \kappa \Ad(\kappa)\,(1 - \cos\theta_j),
\label{eq:kl_vmf_col}
\end{equation}
where $\theta_j \coloneqq \arccos(\hat{\mu}_j^\top \mu_{0,j})$ is the
angle between the posterior and prior mean directions for column $j$.

Combining~\eqref{eq:kl_decomp}, \eqref{eq:kl_gaussian},
and~\eqref{eq:kl_vmf_col} with Theorem~\ref{thm:mcallester}:

\begin{theorem}[PAC-Bayes Bound for DoRA]
\label{thm:pac_dora}
Under our assumptions, let $\theta_j$ denote the angle between
$\mu_{0,j}$ and $\hat{\mu}_j$.
For any $\delta \in (0,1)$, with probability at least $1 - \delta$:
\[
  R(Q) \leq \hat{R}(Q) + \sqrt{
  \frac{\displaystyle\sum_{j=1}^d
    \left[\frac{(\hat{m}_j - m_{0,j})^2}{2\sigma^2}
    + \kappa \Ad(\kappa)(1 - \cos\theta_j)\right]
    + \ln\frac{2\sqrt{n}}{\delta}}{2n}}.
\]
\end{theorem}

When $\kappa \gg d$ and $\theta_j \ll 1$:
\begin{itemize}
  \item $\Ad(\kappa) = 1 - O(d/\kappa)$, so $\kappa\Ad(\kappa) \approx \kappa$.
  \item $1 - \cos\theta_j = \tfrac{\theta_j^2}{2} + O(\theta_j^4)$.
\end{itemize}
These follow from the Bessel function asymptotic $I_\nu(\kappa) \sim
e^\kappa / \sqrt{2\pi\kappa}$ as $\kappa \to \infty$ and from the
standard Taylor expansion of cosine. Therefore,
\begin{equation}
  \KL(Q \| P) \approx \sum_{j=1}^d
  \left[\frac{(\hat{m}_j - m_{0,j})^2}{2\sigma^2}
  + \frac{\kappa}{2}\theta_j^2\right].
\label{eq:kl_highconc_dora}
\end{equation}
This is a sum of squared magnitude deviations and squared angular
deviations, with the angular penalty weighted by $\kappa/2$.

% ---------------------------------------------------------------
\subsection{Global Decomposition: MAP}
\label{sec:map}

In MAP, the parameter vector $\Theta \in \R^k$ (the flattened weight
matrix) is decomposed globally:
\[
  \Theta = \rho\, U, \quad \rho \in \R_+,\; U \in \Sk.
\]
The prior and posterior are:
\[
  P = \mathcal{N}(\rho_0, \sigma^2)\cdot\vMF(U_0, \kappa),
  \qquad
  Q = \mathcal{N}(\hat\rho, \sigma^2)\cdot\vMF(\hat{U}, \kappa).
\]
Since $P$ and $Q$ both factor over magnitude and direction, the KL
decomposes as:
\begin{align}
  \KL(Q \| P)
  &= \KL\bigl(\mathcal{N}(\hat\rho,\sigma^2) \| \mathcal{N}(\rho_0,\sigma^2)\bigr)
   + \KL\bigl(\vMF(\hat{U},\kappa) \| \vMF(U_0,\kappa)\bigr) \notag\\
  &= \frac{(\hat\rho - \rho_0)^2}{2\sigma^2}
   + \kappa \Ak(\kappa)\,(1 - \cos\Theta_{\mathrm{global}}),
\label{eq:kl_map}
\end{align}
where $\Theta_{\mathrm{global}} \coloneqq \arccos(\hat{U}^\top U_0)$ and
$\Ak(\kappa) = I_{k/2}(\kappa)/I_{k/2-1}(\kappa)$.

\begin{theorem}[PAC-Bayes Bound for MAP]
\label{thm:pac_map}
For any $\delta \in (0,1)$, with probability at least $1 - \delta$:
\[
  R(Q) \leq \hat{R}(Q) + \sqrt{
  \frac{\dfrac{(\hat\rho - \rho_0)^2}{2\sigma^2}
    + \kappa \Ak(\kappa)(1 - \cos\Theta_{\mathrm{global}})
    + \ln\frac{2\sqrt{n}}{\delta}}{2n}}.
\]
In the high-concentration regime ($\kappa \gg k$, $\Theta_{\mathrm{global}} \ll 1$):
\[
  \KL(Q \| P) \approx \frac{(\hat\rho - \rho_0)^2}{2\sigma^2}
  + \frac{\kappa}{2}\Theta_{\mathrm{global}}^2.
\]
\end{theorem}

\begin{proof}
Substitute~\eqref{eq:kl_map} into Theorem~\ref{thm:mcallester}.
The high-concentration approximation uses the same Bessel asymptotics
as in the DoRA case, with $d$ replaced by $k$.
\end{proof}

% ---------------------------------------------------------------
\subsection{A Data-Dependent PAC-Bayesian Bound}
\label{sec:catoni}

The McAllester bound scales as $\sqrt{\KL / n}$, which can be
conservative when the KL is small.
We now state a Catoni-style bound~\cite{catoni2007pac} that scales
\emph{linearly} in the KL.

\begin{theorem}[Catoni PAC-Bayes Bound]
\label{thm:catoni}
Let $\phikl(x) = \tfrac{1}{\lambda}\ln(1 + \lambda x)$ and define
the Catoni empirical risk:
\[
  \Rhatl(Q) \coloneqq \E_{h \sim Q}\frac{1}{n}\sum_{i=1}^n
  \phikl\!\bigl(\ell(h, z_i)\bigr).
\]
For any $\lambda > 0$, with probability at least $1 - \delta$ over $S$,
for all posteriors $Q$:
\[
  R(Q) \leq \frac{1}{1 - e^{-\lambda}}
  \left(\Rhatl(Q) + \frac{\KL(Q\|P) + \ln\frac{1}{\delta}}{\lambda n}\right).
\]
\end{theorem}

\begin{proof}[Proof sketch]
The key is that for $x \in [0,1]$, $\phikl(x) = \tfrac{1}{\lambda}\ln(1+\lambda x)
\geq x - \tfrac{\lambda}{2}x^2$.  Applying the Donsker–Varadhan inequality
to the tilted loss $e^{\lambda \ell(h,z_i)}$, one obtains the bound via
the moment generating function of the loss.
Specifically:
\begin{align*}
  \E_{h\sim Q}\E_S\left[\frac{1}{n}\sum_i e^{\lambda \ell(h,z_i)}\right]
  &\leq e^{\KL(Q\|P)} \cdot
    \E_{h\sim P}\E_S\left[\frac{1}{n}\sum_i e^{\lambda \ell(h,z_i)}\right].
\end{align*}
Applying $\ln(1+x) \geq x/(1+x)$ and rearranging yields the stated bound;
see~\cite{catoni2007pac} for the full proof.

%% [EDITORIAL]: Consider including a full proof here (it is short) to
%% make the paper self-contained for a CS 229 audience.
\end{proof}

\begin{remark}
Unlike McAllester's bound, Theorem~\ref{thm:catoni}:
\begin{itemize}
  \item scales linearly in $\KL(Q\|P)$ rather than $\sqrt{\KL}$,
    giving tighter bounds when $\KL$ is small;
  \item allows tuning $\lambda$ to trade off empirical risk bias and
    complexity penalty;
  \item reduces to McAllester's bound (up to constants) when
    $\lambda = O(1/\sqrt{n})$.
\end{itemize}
\end{remark}

\paragraph{Application to weight-decomposed adaptation.}
Substituting the KL bounds from Sections~\ref{sec:dora} and~\ref{sec:map}
into Theorem~\ref{thm:catoni}:
\[
  R(Q) \lesssim \Rhatl(Q) + \frac{1}{\lambda n}
  \left(C_{\mathrm{mag}} + C_{\mathrm{dir}} + \ln\frac{1}{\delta}\right),
\]
where:
\begin{align*}
  C^{\mathrm{DoRA}}_{\mathrm{dir}}
  &= \kappa \Ad(\kappa) \sum_{j=1}^d (1 - \cos\theta_j),\\
  C^{\mathrm{MAP}}_{\mathrm{dir}}
  &= \kappa \Ak(\kappa)(1 - \cos\Theta_{\mathrm{global}}),\\
  C_{\mathrm{mag}}
  &= \sum_{j=1}^d \frac{(\hat{m}_j - m_{0,j})^2}{2\sigma^2}
    \quad \text{(DoRA)}, \quad
    \text{or}\quad
    \frac{(\hat\rho - \rho_0)^2}{2\sigma^2}
    \quad \text{(MAP)}.
\end{align*}

% ---------------------------------------------------------------
\subsection{Comparison Between Local and Global Adaptation}
\label{sec:comparison}

We now compare the directional complexity terms of DoRA and MAP.
The following theorem is the geometric heart of the paper.

\begin{theorem}[Angular Complexity Comparison]
\label{thm:comparison}
Let $u_j, \hat{u}_j \in \Sd$ be the pretrained and fine-tuned unit
direction vectors for column $j$ in DoRA, with $\theta_j \in [0, \pi/2]$
(as holds in fine-tuning).
Define the concatenated unit vectors:
\[
  U \coloneqq \frac{1}{\sqrt{d}}(u_1, \ldots, u_d) \in \mathbb{S}^{dk-1},
  \qquad
  \hat{U} \coloneqq \frac{1}{\sqrt{d}}(\hat{u}_1, \ldots, \hat{u}_d).
\]
Then $\Theta_{\mathrm{global}} \coloneqq \arccos(\hat{U}^\top U)$ satisfies:
\begin{equation}\label{eq:thm_comparison}
  \Theta_{\mathrm{global}}^2
  \;\leq\; \frac{1}{d}\sum_{j=1}^d \theta_j^2
  \;\leq\; \sum_{j=1}^d \theta_j^2.
\end{equation}
Moreover, the following identity holds exactly:
\begin{equation}\label{eq:cos_identity}
  1 - \cos\Theta_{\mathrm{global}}
  = \frac{1}{d}\sum_{j=1}^d (1 - \cos\theta_j).
\end{equation}
\end{theorem}

\begin{proof}
\textbf{Step 1 (exact identity).}
From the definition of $\hat{U}$ and $U$:
\[
  \cos\Theta_{\mathrm{global}} = \hat{U}^\top U
  = \frac{1}{d}\sum_{j=1}^d \hat{u}_j^\top u_j
  = \frac{1}{d}\sum_{j=1}^d \cos\theta_j,
\]
so the following identity holds \emph{exactly}:
\begin{equation}
  1 - \cos\Theta_{\mathrm{global}}
  = \frac{1}{d}\sum_{j=1}^d (1 - \cos\theta_j).
  \label{eq:cos_id_proof}
\end{equation}

\textbf{Step 2 (Jensen's inequality).}
Under the assumption $\theta_j \in [0,\pi/2]$, the function $\cos$ is
concave since $\cos'' x = - \cos x \le 0$ for $x \in [0,\pi/2]$.  Jensen's inequality for convex functions gives:
\[
  \frac{1}{d}\sum_{j=1}^d \cos\theta_j
  \;\geq\; \cos\!\left(\frac{1}{d}\sum_{j=1}^d \theta_j\right)
  = \cos\bar\theta,
  \quad \bar\theta \coloneqq \frac{1}{d}\sum_{j=1}^d \theta_j.
\]
Since $\arccos$ is strictly decreasing:
\[
  \Theta_{\mathrm{global}}
  = \arccos\!\left(\frac{1}{d}\sum_{j=1}^d\cos\theta_j\right)
  \;\leq\; \arccos(\cos\bar\theta)
  = \bar\theta.
\]

\textbf{Step 3 (Cauchy--Schwarz).}
Squaring and applying the QM--AM inequality
$(\frac{1}{d}\sum\theta_j)^2 \leq \frac{1}{d}\sum\theta_j^2$:
\[
  \Theta_{\mathrm{global}}^2
  \;\leq\; \bar\theta^2
  = \left(\frac{1}{d}\sum_{j=1}^d\theta_j\right)^{\!2}
  \;\leq\; \frac{1}{d}\sum_{j=1}^d\theta_j^2
  \;\leq\; \sum_{j=1}^d\theta_j^2.
\]
The stronger bound $\Theta_{\mathrm{global}}^2 \leq \frac{1}{d}\sum_j\theta_j^2$
also holds directly.
\end{proof}

\begin{remark}
Theorem~\ref{thm:comparison} reveals the geometric intuition concisely:
the global angle $\Theta_{\mathrm{global}}$ is an \emph{average} over local
angles (in the sense of averaging their cosines), whereas the DoRA
directional complexity $C^{\mathrm{DoRA}}_{\mathrm{dir}}$ is a \emph{sum}.
Consequently, MAP always has a smaller or equal \emph{global} angular
deviation, but this comes at the cost of losing per-column resolution.
\end{remark}

\paragraph{Implications.}
As a result:
\begin{itemize}
  \item MAP incurs a smaller directional KL penalty when updates are
    highly aligned across columns.
  \item DoRA achieves tighter bounds when updates are sparse or
    localized (many $\theta_j = 0$).
\end{itemize}
This highlights a fundamental tradeoff between expressivity and
generalization complexity.

% ---------------------------------------------------------------
\subsection{When Does DoRA Achieve a Tighter Bound?}
\label{sec:dora_advantage}

\begin{theorem}[DoRA Advantage Under Localized Updates]
\label{thm:dora_advantage}
Suppose the update is sparse: for some set $\mathcal{S} \subseteq \{1,\ldots,d\}$
with $|\mathcal{S}| = s \ll d$, $\theta_j = 0$ for $j \notin \mathcal{S}$ and
$\theta_j \leq \varepsilon$ for $j \in \mathcal{S}$.
Then:
\[
  C^{\mathrm{DoRA}}_{\mathrm{dir}} \leq \kappa \Ad(\kappa) \cdot s\,(1 - \cos\varepsilon).
\]
Moreover, under the exact identity~\eqref{eq:cos_identity} of
Theorem~\ref{thm:comparison}, and using $\theta_j = 0$ for $j \notin
\mathcal{S}$:
\begin{equation}
  1 - \cos\Theta_{\mathrm{global}}
  = \frac{1}{d}\sum_{j \in \mathcal{S}} (1-\cos\theta_j)
  \leq \frac{s}{d}(1 - \cos\varepsilon),
\label{eq:global_upper}
\end{equation}
where the inequality uses $1-\cos\theta_j \leq 1-\cos\varepsilon$ for
$\theta_j \leq \varepsilon$ (since $1-\cos$ is increasing on $[0,\pi]$).
Therefore:
\[
  C^{\mathrm{MAP}}_{\mathrm{dir}} \leq \kappa \Ak(\kappa)\cdot\frac{s}{d}(1 - \cos\varepsilon).
\]
\begin{remark}
Both DoRA and MAP thus admit \emph{upper} bounds proportional to $s(1-\cos\varepsilon)$
in the sparse regime.
Comparing the two upper bounds:
\[
  \frac{C^{\mathrm{MAP}}_{\mathrm{dir}}\text{ upper bound}}
       {C^{\mathrm{DoRA}}_{\mathrm{dir}}\text{ upper bound}}
  = \frac{\Ak(\kappa)/d}{\Ad(\kappa)}
  = \frac{\Ak(\kappa)}{d\,\Ad(\kappa)}.
\]
Since $\Ak(\kappa) < \Ad(\kappa)$ for $k > d$ (as $A$ is decreasing in dimension),
the MAP upper bound is always smaller by a factor $\Ak(\kappa)/(d\,\Ad(\kappa)) < 1/d$.
Using the leading-order approximation $A_n(\kappa) \approx \kappa/n$ for $n \gg \kappa$:
this ratio is approximately $\kappa/k$ (the inverse full dimension).

However, this does \emph{not} mean MAP is universally preferred.
The comparison is between priors on spheres of very different dimensions ($d$
vs.\ $k = d^2$ for a square weight matrix).
With fixed $\kappa$, the MAP prior on $\mathbb{S}^{k-1}$ is nearly flat
($\Ak(\kappa) \approx \kappa/k \approx 0$), providing essentially no
directional regularization for large $k$.
DoRA's per-column prior, by contrast, remains informative ($\Ad(\kappa) \approx
\kappa/d$).
The correct interpretation is that DoRA provides a \emph{more discriminating}
complexity measure: its bound distinguishes between a model where $s$ specific
columns rotated (sparse) vs.\ all $d$ columns rotating equally (dense),
whereas MAP conflates both into a single global angle with negligible KL.
For a fair quantitative comparison, $\kappa$ should be scaled by the
respective sphere dimension (e.g.\ $\kappa_{\mathrm{eff}} = \kappa \cdot d$ for
DoRA and $\kappa_{\mathrm{eff}} = \kappa \cdot k$ for MAP), a calibration
question addressed in Section~\ref{sec:disc_kappa}.
\end{remark}
\end{theorem}

\begin{proof}
The upper bound on $C^{\mathrm{DoRA}}_{\mathrm{dir}}$ follows directly: since
$\theta_j = 0$ for $j \notin \mathcal{S}$, those terms contribute zero to
the sum.  For $j \in \mathcal{S}$, $(1-\cos\theta_j) \leq (1-\cos\varepsilon)$.

For the MAP upper bound: equation~\eqref{eq:global_upper} gives
$1 - \cos\Theta_{\mathrm{global}} \leq \tfrac{s}{d}(1-\cos\varepsilon)$
(since $\theta_j = 0$ for $j \notin \mathcal{S}$ and $1-\cos\theta_j \leq 1-\cos\varepsilon$
for $\theta_j \leq \varepsilon$).
The MAP directional complexity bound then follows by substituting into the
upper bound of Theorem~\ref{thm:nonasymptotic}.
\end{proof}

\begin{remark}[On the DoRA vs.\ MAP comparison]
\label{rem:comparison}
The ratio of upper bounds from Theorem~\ref{thm:dora_advantage} is
$d\,\Ad(\kappa)/\Ak(\kappa)$.
Since $A$ is \emph{decreasing} in dimension, $\Ak(\kappa) < \Ad(\kappa)$ for
$k > d$, so this ratio exceeds $d > 1$: DoRA's upper bound is always larger
than MAP's.

The condition $d\,\Ad(\kappa) < \Ak(\kappa)$ stated in an earlier version of
this corollary is \emph{never satisfied} in practice (it would require $A$ to
be increasing, contradicting the Bessel ratio monotonicity).
The correct conclusion is that the two methods trade off \emph{discriminability}
vs.\ \emph{raw KL magnitude}: DoRA resolves which columns changed;
MAP achieves a smaller KL by using a diffuse prior on the global direction.
The practical guidance in Table~\ref{tab:guidance} rests on discriminability
(which bound best reflects the actual update structure), not on which
literal KL value is smaller.
\end{remark}

\paragraph{Interpretation.}
DoRA achieves a strictly tighter bound when:
\begin{itemize}
  \item updates are sparse (small $s/d$),
  \item directional changes are localized to few columns,
  \item global reorientation would unnecessarily rotate unaffected components.
\end{itemize}

\paragraph{Dense regime.}
Conversely, when all columns rotate coherently ($\theta_j \approx
\Theta_{\mathrm{global}}$ for all $j$), the per-column angles sum to
$\sum_j(1-\cos\theta_j) \approx d(1-\cos\Theta_{\mathrm{global}})$,
giving $C^{\mathrm{DoRA}}_{\mathrm{dir}} \approx d\cdot C^{\mathrm{MAP}}_{\mathrm{dir}}$.
Thus $C^{\mathrm{MAP}}_{\mathrm{dir}} \ll C^{\mathrm{DoRA}}_{\mathrm{dir}}$ in this regime,
and MAP is preferred.

%% [EDITORIAL]: This is the most compelling theoretical contribution.
%% Consider adding a simple synthetic experiment (e.g., a
%% $2\times 2$ matrix) illustrating the two regimes with concrete numbers.

% ============================================================
\section{Data-Dependent Posteriors via Gibbs Measures}
\label{sec:gibbs}

%% [EDITORIAL]: Section 2.5 of the original draft placed the Catoni
%% bound \emph{before} the comparison theorem; the narrative flow is
%% better with the comparison first (so the reader understands what is
%% being optimized) and the Gibbs section second.
%% The current restructuring follows this logic.

\subsection{Gibbs Posterior}
\label{sec:gibbs_def}

Given the Catoni objective:
\[
  \mathcal{L}_\lambda(Q) \coloneqq \Rhatl(Q) + \frac{1}{\lambda n}\KL(Q \| P),
\]
the unique minimizer over all distributions $Q$ is the \emph{Gibbs posterior}:
\begin{equation}
  \frac{dQ^*(h)}{dP(h)}
  = \frac{\exp\bigl(-\lambda n\,\Rhatl(h)\bigr)}{Z_\lambda},
  \qquad
  Z_\lambda = \int \exp\bigl(-\lambda n\,\Rhatl(h)\bigr)\,dP(h).
\label{eq:gibbs}
\end{equation}
The Gibbs posterior balances data fit (via $\Rhatl(h)$) and proximity to
the prior (via the implicit $\KL$ regularization).

\subsection{Gibbs Posterior for Weight-Decomposed Adaptation}
\label{sec:gibbs_wda}

Under our assumptions and the structure of DoRA, the
Gibbs posterior is:
\[
  Q^*\!\bigl(\{m_j, u_j\}_{j=1}^d\bigr)
  \propto \prod_{j=1}^d P_{\mathrm{mag},j}(m_j)\,P_{\mathrm{dir},j}(u_j)
  \cdot\exp\bigl(-\lambda n\,\Rhatl(\{m_j u_j\})\bigr).
\]
While the loss couples all columns through the forward pass, we obtain a
tractable approximation via mean-field:
\[
  Q(\{m_j, u_j\}) = \prod_{j=1}^d Q_{\mathrm{mag},j}(m_j)\,Q_{\mathrm{dir},j}(u_j).
\]
Minimizing $\E_Q[\Rhatl(h)] + \tfrac{1}{\lambda n}\KL(Q\|P)$ over
mean-field families and substituting the KL formulas from
section 3.3:

\begin{equation}
  \mathcal{L}_{\mathrm{DoRA}}
  = \E_Q[\Rhatl(h)]
  + \frac{1}{\lambda n}\sum_{j=1}^d
    \left[\frac{(\hat{m}_j - m_{0,j})^2}{2\sigma^2}
    + \kappa \Ad(\kappa)(1-\cos\theta_j)\right].
\label{eq:obj_dora}
\end{equation}
Similarly, for MAP:
\begin{equation}
  \mathcal{L}_{\mathrm{MAP}}
  = \E_Q[\Rhatl(h)]
  + \frac{1}{\lambda n}
    \left[\frac{(\hadt\rho - \rho_0)^2}{2\sigma^2}
    + \kappa \Ak(\kappa)(1-\cos\Theta_{\mathrm{global}})\right].
\label{eq:obj_map}
\end{equation}

\paragraph{Connection to existing regularization.}
The objectives \eqref{eq:obj_dora} and \eqref{eq:obj_map} each have two
parts:
\begin{itemize}
  \item \textbf{Weight decay}: the $(\hat{m}-m_0)^2/2\sigma^2$ term
    penalizes deviation of magnitudes from the pretrained values.
  \item \textbf{Cosine regularization}: the $\kappa\Ad(\kappa)(1-\cos\theta)$
    term penalizes angular deviation from the pretrained directions.
\end{itemize}
These penalties arise from first principles via the PAC-Bayesian framework
rather than being introduced ad hoc—providing a theoretical justification
for cosine-based regularizers used empirically in PEFT.

%% [EDITORIAL]: This connection is one of the most elegant contributions.
%% Consider adding a forward-pointer to how $\kappa$ acts as the
%% regularization coefficient: larger $\kappa$ means stronger directional
%% regularization.  Practitioners can tune $\kappa$ as a hyperparameter
%% and interpret it as prior confidence in the pretrained directions.

% ============================================================
\section{Theoretical Predictions and Planned Experiments}
\label{sec:experiments}

\noindent\textbf{Scope of this section.}
The four predictions stated below follow directly from the theoretical
analysis of Sections~\ref{sec:pac_analysis}--\ref{sec:comparison}.
This section states each prediction in testable, falsifiable form,
specifies the experimental protocol required to evaluate it, and
identifies what confirming or disconfirming evidence would look like.
\emph{No experiments have been conducted; this section constitutes a
pre-registered plan for future empirical work.}

% ---------------------------------------------------------------
\subsection{Experimental Setup}
\label{sec:exp_setup}

\paragraph{Models and tasks.}
We will fine-tune RoBERTa-base (125\,M parameters)~\cite{liu2019roberta}
on six GLUE tasks (SST-2, MNLI, QQP, QNLI, MRPC, RTE) and optionally
LLaMA-2-7B~\cite{touvron2023llama} on instruction-following using the
Alpaca-52k dataset~\cite{alpaca2023}.
We will compare three PEFT methods --- LoRA (rank $r = 8$), DoRA
(rank $r = 8$), and MAP (rank $r = 8$) --- matched for total trainable
parameter count.
All runs will use AdamW with a cosine learning-rate schedule, peak
learning rate $2 \times 10^{-4}$, and three random seeds; we will report
mean and standard deviation across seeds.

\paragraph{Computing $\theta_j$ and $\Theta_{\mathrm{global}}$.}
After fine-tuning, for each weight matrix $W \in \R^{d_{\mathrm{out}}
\times d_{\mathrm{in}}}$ we will compute the column-wise angle
\[
  \theta_j = \arccos\!\left(
    \frac{w_{0,j}^\top \hat{w}_j}{\|w_{0,j}\|\,\|\hat{w}_j\|}
  \right), \quad j = 1, \ldots, d_{\mathrm{in}},
\]
where $w_{0,j}$ and $\hat{w}_j$ are the $j$-th columns of the
pretrained and fine-tuned weight matrices, respectively.
The global angle is computed on the flattened vectors:
$\Theta_{\mathrm{global}} = \arccos(\hat{w}^\top w_0 /
(\|\hat{w}\|\,\|w_0\|))$ where $w_0, \hat{w} \in \R^{d_{\mathrm{out}}
d_{\mathrm{in}}}$.

\paragraph{Directional complexity estimator.}
We will estimate:
\[
  \widehat{C}^{\mathrm{DoRA}}_{\mathrm{dir}}
    = \kappa \sum_{j=1}^{d_{\mathrm{in}}} (1 - \cos\theta_j),
  \qquad
  \widehat{C}^{\mathrm{MAP}}_{\mathrm{dir}}
    = \kappa\,(1 - \cos\Theta_{\mathrm{global}}),
\]
using $\kappa = 10$ as the default, with an ablation over
$\kappa \in \{1, 5, 10, 50, 100\}$ (Section~\ref{sec:ablation}).
The effective sparsity is $s = |\{j : \theta_j > \tau\}|$,
$\tau = 0.01$ radians.
Layer-level estimates will be aggregated over all weight matrices of the
same type (attention vs.\ MLP) and averaged across layers.

% ---------------------------------------------------------------
\subsection{Prediction P1: $\widehat{C}_{\mathrm{dir}}$ Predicts Test Error}
\label{sec:kl_predict}

\paragraph{Theoretical basis.}
The PAC-Bayes bound (Theorem~\ref{thm:mcallester}) implies
\[
  \mathrm{err}(h) \;\lesssim\; \hat{R}_n(h) + \sqrt{\frac{C_{\mathrm{dir}}}{n}},
\]
so $C_{\mathrm{dir}}$ should be positively correlated with the
generalisation gap, and hence with test error after controlling for
training error.

\paragraph{Protocol.}
For each (method, task, layer-type, seed) combination we will record the
pair $(\widehat{C}_{\mathrm{dir}}, \mathrm{err})$ where $\mathrm{err}$ is
the task-specific test error (1 $-$ accuracy for classification; perplexity
normalised to $[0,1]$ via $1 - e^{-\Delta}$ for generation).
We will pool all 108 pairs
(3 methods $\times$ 6 tasks $\times$ 2 layer types $\times$ 3 seeds)
and compute:
\begin{itemize}
  \item \textbf{Spearman rank correlation} $\rho_S$ between
    $\widehat{C}_{\mathrm{dir}}$ and $\mathrm{err}$.
  \item \textbf{Pearson correlation} $\rho_P$ between
    $\ln\widehat{C}_{\mathrm{dir}}$ and $\mathrm{err}$, testing the
    log-linear relationship implied by the $\sqrt{C/n}$ scaling.
\end{itemize}
Statistical significance will be assessed via a permutation test with
$B = 10{,}000$ bootstrap resamples; we will report $p$-values and
95\% confidence intervals.

\paragraph{Expected outcome.}
We predict $\rho_S > 0.5$ with $p < 0.01$, with the log-scale Pearson
correlation providing stronger fit than the linear-scale correlation,
consistent with the $\sqrt{C_{\mathrm{dir}}/n}$ scaling of the bound.
A scatter plot of $\widehat{C}_{\mathrm{dir}}$ (log scale) vs.\ test error
will visualise the relationship.

\paragraph{Falsification criterion.}
A result $\rho_S < 0.2$ or $p > 0.05$ would call into question whether
the angular complexity measure captures the dominant generalisation-gap
driver in practice.

% ---------------------------------------------------------------
\subsection{Prediction P2: Sparse vs.\ Dense Regime Crossover}
\label{sec:sparse_dense}

\paragraph{Theoretical basis.}
Theorem~\ref{thm:dora_advantage} establishes that
$C^{\mathrm{DoRA}}_{\mathrm{dir}} < C^{\mathrm{MAP}}_{\mathrm{dir}}$
when the update is sparse ($s/d$ small), with MAP preferred when updates
are globally aligned ($s/d$ large).
The crossover occurs near $s/d \approx 0.5$ under uniform-angle
assumptions.

\paragraph{Protocol.}
We will vary the sparsity ratio $s/d$ by sweeping rank
$r \in \{2, 4, 8, 16, 32\}$ in DoRA and MAP, with all other
hyperparameters fixed.
Lower rank induces more localised updates (smaller $s/d$).
We will plot $\widehat{C}_{\mathrm{dir}}$ for DoRA and MAP as a function
of $s/d$ on SST-2 (representative task) and record the empirical crossover
rank.

\paragraph{Expected outcome.}
DoRA should achieve strictly lower $\widehat{C}_{\mathrm{dir}}$ at
$s/d \lesssim 0.3$, with MAP preferred at $s/d \gtrsim 0.7$, and the
curves crossing near $s/d \approx 0.5$.

\paragraph{Falsification criterion.}
Absence of a crossover (DoRA always preferred, or MAP always preferred)
would contradict Theorem~\ref{thm:dora_advantage} and suggest the
sparsity-ratio parameterisation does not match the theoretical $s/d$ regime.

% ---------------------------------------------------------------
\subsection{Prediction P3: Attention vs.\ MLP Angular Variance}
\label{sec:alignment}

\paragraph{Theoretical basis.}
The layer-specific guidance in Table~\ref{tab:guidance} rests on the
empirical finding of Tanwar et al.~\cite{tanwar2025domain} that MLP
layers undergo larger directional updates during fine-tuning than
attention layers.
Our framework translates this into a prediction about
$\mathrm{Var}(\theta_j)$: if MLP layers consistently show higher angular
variance, MAP is appropriate for attention and DoRA is appropriate for MLP.

\paragraph{Protocol.}
For each fine-tuned model we will compute $\mathrm{Var}(\theta_j)$
separately over attention weight matrices ($W_Q, W_K, W_V, W_O$) and
MLP weight matrices ($W_{\mathrm{up}}, W_{\mathrm{gate}},
W_{\mathrm{down}}$), averaged across layers and seeds.
Table~\ref{tab:variance} will report the MLP/attention ratio for each
task upon completion of the experiments.

\begin{table}[ht]
\centering
\caption{Planned measurement: mean per-column angular variance
  $\mathrm{Var}(\theta_j)$ (rad$^2$) for attention vs.\ MLP layers,
  averaged over GLUE and Alpaca tasks.
  Values to be filled upon experiment completion.
  The theoretical prediction is MLP/Attn ratio $> 1$ across all tasks,
  consistent with \cite{tanwar2025domain}.}
\label{tab:variance}
\begin{tabular}{lccc}
\toprule
Task & Attn $\mathrm{Var}(\theta_j)$ &
  MLP $\mathrm{Var}(\theta_j)$ & Ratio (MLP/Attn) \\
\midrule
SST-2  & --- & --- & (predicted $> 1$) \\
MNLI   & --- & --- & (predicted $> 1$) \\
QQP    & --- & --- & (predicted $> 1$) \\
QNLI   & --- & --- & (predicted $> 1$) \\
MRPC   & --- & --- & (predicted $> 1$) \\
RTE    & --- & --- & (predicted $> 1$) \\
\midrule
Mean   & TBD & TBD & TBD \\
\bottomrule
\end{tabular}
\end{table}

\paragraph{Falsification criterion.}
If attention and MLP layers show comparable angular variance (ratio $< 1.5$),
the layer-specific guidance in Table~\ref{tab:guidance} should be revised
to treat all layers uniformly.

% ---------------------------------------------------------------
\subsection{Prediction P4: Non-Asymptotic vs.\ Asymptotic Bound Tightness}
\label{sec:ablation}

\paragraph{Theoretical basis.}
The asymptotic approximation replaces $\Ad(\kappa)$ with 1, overestimating
the true directional KL by a factor of $1/\Ad(\kappa)$.
For moderate $\kappa$ and large $d$ (as in transformer weight matrices),
$\Ad(\kappa) \ll 1$, so the non-asymptotic bound
(Theorem~\ref{thm:nonasymptotic}) is strictly tighter.
The exact values of $\Ad(\kappa)$ for $d = 768$ (RoBERTa-base hidden
dimension) can be computed numerically as
$\Ad(\kappa) = I_{d/2}(\kappa)/I_{d/2-1}(\kappa)$;
for $d = 768$ and $\kappa \leq 100$, $\Ad(\kappa) \ll 1$, so the
non-asymptotic correction is large.

\paragraph{Protocol.}
For $\kappa \in \{1, 5, 10, 50, 100\}$ we will compute both the
asymptotic bound ($\kappa\sum(1-\cos\theta_j)$) and the non-asymptotic
bound ($\kappa\Ad(\kappa)\sum(1-\cos\theta_j)$) on the empirical
$\{\theta_j\}$ values from RoBERTa-base fine-tuned on SST-2
(DoRA, seed 1).
We will report the overestimate ratio and compare the corresponding
PAC-Bayes bound values.

\paragraph{Expected outcome.}
For $\kappa \leq 10$, the overestimate should be substantial (tens of
percent to orders of magnitude), shrinking toward zero as $\kappa \to \infty$.
The Catoni $\lambda^*$ ablation (Theorem~\ref{thm:catoni}) should yield
a materially tighter bound than the McAllester variant
(Theorem~\ref{thm:mcallester}) at the optimal $\lambda^*$.

\paragraph{Falsification criterion.}
If $\Ad(\kappa) \approx 1$ for the $\kappa$ values used in practice,
the non-asymptotic correction would be negligible and the asymptotic
formula would be adequate.
Given that $d = 768$, this outcome is highly unlikely but would imply
that either $\kappa \gg d$ in typical fine-tuning or that the
column-wise vMF model is miscalibrated.

% ---------------------------------------------------------------
\subsection{Summary of Predictions}

Table~\ref{tab:exp_summary} summarises the four predictions and the
criteria by which they will be evaluated.

\begin{table}[ht]
\centering
\caption{Theoretical predictions and planned validation criteria.}
\label{tab:exp_summary}
\begin{tabular}{p{3.5cm}p{3cm}p{4.5cm}}
\toprule
Prediction & Key metric & Success criterion \\
\midrule
(P1) $C_{\mathrm{dir}}$ predicts test error &
  Spearman $\rho_S$ &
  $\rho_S > 0.5$, $p < 0.01$, pooled over 108 data points \\
(P2) DoRA preferred in sparse regime &
  Crossover $s/d$ &
  DoRA $<$ MAP at $s/d \leq 0.3$; MAP $<$ DoRA at $s/d \geq 0.7$ \\
(P3) MLP Var$(\theta_j) >$ Attn &
  Variance ratio &
  MLP/Attn ratio $> 1.5$ across all GLUE tasks \\
(P4) Non-asym.\ bound tighter &
  Overestimate ratio at $\kappa=10$ &
  Overestimate $> 10\%$ for $d=768$ \\
\bottomrule
\end{tabular}
\end{table}

% ============================================================
% ============================================================

\section{Discussion}
\label{sec:discussion}

\subsection{Practical Guidance: When to Use DoRA vs.\ MAP}
\label{sec:disc_guidance}

Our theoretical results yield directly actionable decision criteria.
The key observable is the \emph{effective sparsity ratio} $s/d$, where
$s = |\{j : \theta_j > \tau\}|$ counts the columns whose pretrained
directions change significantly (beyond a threshold $\tau$), out of $d$
total columns.
Theorem~\ref{thm:dora_advantage} and Remark~\ref{rem:comparison} together imply:

\begin{itemize}
  \item \textbf{Use DoRA when $s/d$ is small.}  When only a sparse subset
    of columns undergo meaningful directional change, column-wise
    decomposition concentrates complexity where the update actually occurs.
    DoRA's directional KL bound scales as $O(s)$ and explicitly tracks
    \emph{which} $s$ columns changed.
    While the MAP upper bound is numerically smaller (Remark~\ref{rem:comparison}),
    this is because the MAP prior on $\mathbb{S}^{k-1}$ (with $k = d^2$) is
    nearly flat and provides negligible directional regularization for large $k$.
    DoRA's bound is the more \emph{discriminating} complexity measure:
    it distinguishes a sparse, localized update from a globally diffuse one,
    which MAP cannot.

  \item \textbf{Use MAP when $s/d \approx 1$ and variance of
    $\{\theta_j\}$ is low.}  When nearly all columns rotate by
    approximately the same angle ($\theta_j \approx \varepsilon$ for
    all $j$), the global rotation captures the shared structure in a single
    degree of freedom, giving
    $C^{\mathrm{MAP}}_{\mathrm{dir}} = O(\kappa(1 - \cos\varepsilon))$
    versus
    $C^{\mathrm{DoRA}}_{\mathrm{dir}} = O(\kappa d (1 - \cos\varepsilon))$,
    a factor of $d$ improvement for MAP.

  \item \textbf{Use variance of $\{\theta_j\}$ as a diagnostic.}
    In practice, the $\theta_j$ can be computed cheaply after a short
    fine-tuning run (or even estimated from gradient signals before
    training).  High variance indicates localized updates and favors
    DoRA; near-zero variance indicates globally aligned updates and
    favors MAP.
\end{itemize}

\paragraph{Layer-specific guidance.}
A practically important refinement emerges from the empirical finding of
Tanwar et al.~\cite{tanwar2025domain} that attention and MLP layers
exhibit qualitatively different directional update patterns during
domain-specific fine-tuning: attention weight updates primarily
\emph{amplify existing directions} in the pretrained subspace, while MLP
updates primarily \emph{inject novel directions} orthogonal to the
pretrained subspace.
In the language of our framework, this means:
\begin{itemize}
  \item \textbf{Attention weight matrices} ($W_Q, W_K, W_V, W_O$) tend to
    have low $\mathrm{Var}(\theta_j)$ and high $s/d$ (many columns rotate
    by small, similar angles) $\Rightarrow$ \textbf{MAP is preferred}:
    the global rotation captures the shared amplification in a single
    angular parameter.
  \item \textbf{MLP weight matrices} ($W_{\mathrm{up}},
    W_{\mathrm{gate}}, W_{\mathrm{down}}$) tend to have high
    $\mathrm{Var}(\theta_j)$ and moderate $s/d$ (a sparse subset of
    columns rotate substantially to inject new information)
    $\Rightarrow$ \textbf{DoRA is preferred}: per-column KL accounts
    for exactly which columns changed and by how much.
\end{itemize}

Relatedly, AdaLoRA~\cite{zhang2023adalora} finds empirically that attention
and MLP modules merit different rank budgets across layers, a finding
consistent with our analysis: modules with sparse, localized updates
benefit from column-wise complexity accounting (DoRA), while modules
with globally uniform updates are better captured by a single rotation
(MAP).

\noindent Table~\ref{tab:guidance} consolidates these recommendations
along two axes---layer type and task proximity---into a single reference
for practitioners.

\begin{table}[ht]
\centering
\caption{Practical guidance for choosing DoRA vs.\ MAP, organized by
  observable regime and transformer layer type.
  Observables $s/d$ and $\mathrm{Var}(\theta_j)$ can be estimated after
  50--100 gradient steps without completing full fine-tuning.
  Empirical support for the attention/MLP distinction comes from
  \cite{tanwar2025domain}; task-proximity effects are established
  in \cite{vu2020exploring,yang2024unveiling}.}
\label{tab:guidance}
\renewcommand{\arraystretch}{1.4}
\setlength{\tabcolsep}{5pt}
\begin{tabular}{p{3.2cm}p{1.5cm}p{2.0cm}p{1.8cm}p{4.5cm}}
\toprule
\textbf{Setting} & \textbf{Typical $s/d$} &
  \textbf{Typical $\mathrm{Var}(\theta_j)$} &
  \textbf{Rec.} & \textbf{Rationale} \\
\midrule
Attention layers \newline ($W_Q, W_K, W_V, W_O$) &
  High (${\approx}1$) & Low &
  MAP & Globally aligned amplification of pretrained directions;
        single global angle captures all column rotations compactly \\
\addlinespace
MLP layers \newline ($W_{\mathrm{up}}, W_{\mathrm{gate}}, W_{\mathrm{down}}$) &
  Low--med. & High &
  DoRA & Sparse injection of new directions; per-column KL is
         $O(s)$ vs.\ MAP's $O(d)$ \\
\addlinespace
Any layer, \newline in-domain task &
  Any & Low &
  MAP \newline (large $\kappa$) & Pretrained directions encode strong
    inductive biases; tight vMF prior prevents unnecessary drift \\
\addlinespace
Any layer, \newline domain-shift task &
  Low--med. & High &
  DoRA \newline (small $\kappa$) & Novel orthogonal directions required;
    column-wise flexibility with diffuse prior allows free rotation \\
\bottomrule
\end{tabular}
\end{table}

\paragraph{Connection to task structure.}
The layer-type patterns in Table~\ref{tab:guidance} align with
empirically observed task-specific fine-tuning dynamics.
Prior work has shown that in-domain fine-tuning tends to produce small,
globally coherent weight updates, while domain-shift tasks produce more
heterogeneous changes concentrated in specific sublayers and weight
subspaces~\cite{vu2020exploring,yang2024unveiling}.
The table synthesizes these two orthogonal axes---task proximity and
layer type---into a joint criterion that practitioners can evaluate
with a short diagnostic run.

\subsection{Calibrating the Concentration Parameter $\kappa$}
\label{sec:disc_kappa}

The concentration $\kappa$ of the vMF prior has a clear interpretation:
it encodes the practitioner's confidence that the pretrained directions
are meaningful for the target task.
Formally, from~\eqref{eq:obj_dora}, the directional penalty is
$\kappa A_d(\kappa)(1 - \cos\theta_j)$, which scales with $\kappa$ and
thus controls how strongly the learned directions are regularized toward
the pretrained directions.

\begin{itemize}
  \item \textbf{Large $\kappa$ (high confidence in pretrained directions).}
    Appropriate when the fine-tuning task is closely related to the
    pretraining distribution---for example, adapting a language model
    pretrained on Wikipedia to a GLUE benchmark task.
    In this regime, the pretrained directions encode strong inductive
    biases~\cite{vu2020exploring}, and large $\kappa$ prevents the adapter
    from drifting away needlessly.
    The regularization is analogous to the strong cosine similarity
    encouragement used in the original DoRA formulation~\cite{liu2024dora}.

  \item \textbf{Small $\kappa$ (diffuse prior, large domain shift).}
    Appropriate when the target task differs substantially from
    pretraining---for example, instruction-tuning a base language model
    or adapting to a specialized domain such as code or biomedicine.
    Here, a small $\kappa$ allows the adapter to rotate freely, reducing
    the regularization bias at the cost of a potentially looser bound.
    This is consistent with observations that domain-shift fine-tuning
    requires larger weight changes~\cite{tanwar2025domain} and that
    applying strong proximity constraints to pretrained weights can
    hurt performance on out-of-domain tasks.

  \item \textbf{$\kappa$ as a hyperparameter.}
    In practice, $\kappa$ can be selected via a validation-set sweep.
    The PAC-Bayes bound provides a computable oracle: the value of
    $\kappa$ that minimizes the right-hand side of
    Theorem~\ref{thm:pac_dora} (or~\ref{thm:pac_map}) on a held-out set
    is the theoretically optimal choice.
    This gives a principled alternative to tuning $\kappa$ by grid search.
\end{itemize}

\paragraph{Relationship to $\kappa$ and SAM.}
There is a suggestive connection between our directional regularizer and
Sharpness-Aware Minimization (SAM)~\cite{bahri2022sharpness,foret2021sam}.
SAM penalizes parameter perturbations in a Euclidean ball of radius $\rho$,
seeking flat minima in the weight space.
The vMF prior with large $\kappa$ similarly penalizes large angular
deviations from the pretrained initialization, but does so
\emph{geometrically on the hypersphere} rather than in the ambient
Euclidean space.
This means our regularizer is rotation-invariant with respect to
rescaling of the weight vector---a property absent from SAM's Euclidean
ball constraint.
Formalizing this connection and understanding whether angular flatness
implies Euclidean flatness (or vice versa) is a promising direction for
future work.

\subsection{Limitations and Extensions}
\label{sec:disc_limitations}

\paragraph{Mean-field assumption (A1).}
Our KL decomposition in~\eqref{eq:kl_decomp} relies on the product
structure $Q = \prod_j Q_j$, which ignores correlations between columns.
In practice, the columns of a weight matrix interact through the forward
pass and are updated jointly by a shared optimizer.
Relaxing this assumption would require a full matrix-variate distribution
on the joint $(u_1, \ldots, u_d)$, for which closed-form KL computations
are substantially harder.
A tractable middle ground is a \emph{block-diagonal} posterior that
clusters correlated columns together; we leave this to future work.

\paragraph{Equal-concentration assumption (A4).}
Fixing $\kappa$ to be the same for all columns simplifies the KL but
ignores the empirically observed heterogeneity of fine-tuning dynamics
across layers and weight types~\cite{zhang2023adalora}.
Extending to per-column concentrations $\kappa_j$ is straightforward
using Appendix~\ref{app:unequal_kappa}: the KL becomes
$\sum_j \KL(\vMF(\hat{\mu}_j, \kappa_{Q,j}) \| \vMF(\mu_{0,j},
\kappa_{P,j}))$, where each term has the form given in
Appendix~\ref{app:unequal_kappa}.
Learning $\kappa_j$ via PAC-Bayes optimization would allow the prior
to automatically allocate directional complexity budget across columns,
analogous to how AdaLoRA allocates rank budget~\cite{zhang2023adalora}.

\paragraph{Non-vacuous numerical bounds.}
A practical limitation shared with most PAC-Bayes analyses is that the
bounds as stated may be numerically vacuous (i.e., greater than 1) for
large models.
Computing non-vacuous bounds requires careful estimation of the partition
function $Z_\lambda$ in~\eqref{eq:gibbs} via Monte Carlo, as demonstrated
for image classifiers in~\cite{dziugaite2017computing}.
Extending such computations to PEFT at scale is a concrete open problem.

\paragraph{Extension to multi-layer interactions.}
Our analysis treats each weight matrix independently.
In deep networks, the composition of layer-wise updates introduces
additional generalization complexity not captured by per-layer bounds.
Compositional PAC-Bayes bounds~\cite{dziugaite2017computing} could
potentially be applied, but the hyperspherical geometry makes this
non-trivial.

% ============================================================
\section{Conclusion}
\label{sec:conclusion}

We developed a PAC-Bayesian framework for weight-decomposed parameter-efficient
fine-tuning, giving the first principled comparison between column-wise
(DoRA) and global (MAP) adaptation from a generalization perspective.
Our main theoretical result---the Angular Complexity Comparison
(Theorem~\ref{thm:comparison}) and the DoRA Advantage theorem
(Theorem~\ref{thm:dora_advantage})---shows that the two methods occupy
complementary regimes: DoRA excels at sparse, localized updates while
MAP is preferred for globally aligned transformations.
The Gibbs posterior derivation (Section~\ref{sec:gibbs}) provides a
principled justification for cosine regularization and weight decay in
PEFT, and the non-asymptotic vMF KL bounds make the theory valid beyond
the high-concentration limit.

%% [EDITORIAL]: Add 1--2 sentences on future directions:
%% (1) Extending to heterogeneous concentration parameters $\kappa_j$
%%     per layer or per column.
%% (2) Non-vacuous numerical bound computation via the PAC-Bayes--Catoni
%%     objective and Monte Carlo estimation of $Z_\lambda$.
%% (3) Connecting the angular complexity measure to sharpness-aware
%%     minimization (SAM) and flat-minima perspectives.

% ============================================================
\begin{thebibliography}{99}
\bibitem{liu2024dora}
S.-Y.\ Liu, C.-Y.\ Wang, H.\ Yin, P.\ Molchanov, Y.-C.\ F.\ Wang,
K.-T.\ Cheng, and M.-H.\ Chen,
``DoRA: Weight-Decomposed Low-Rank Adaptation,''
\textit{arXiv:2402.09353}, 2024.

\bibitem{si2025map}
C.\ Si, Z.\ Shi, Y.\ Wang, X.\ Yang, S.\ Rahardja, and W.\ Shen,
``MAP: Revisiting Weight Decomposition for Low-Rank Adaptation,''
\textit{arXiv:2505.23094}, 2025.

\bibitem{hu2021lora}
E.\ J.\ Hu, Y.\ Shen, P.\ Wallis, Z.\ Allen-Zhu, Y.\ Li, S.\ Wang,
L.\ Wang, and W.\ Chen,
``LoRA: Low-Rank Adaptation of Large Language Models,''
\textit{arXiv:2106.09685}, 2021.

\bibitem{mcallester1999pac}
D.\ A.\ McAllester,
``Some PAC-Bayesian Theorems,''
\textit{Mach.\ Learn.}, 37(3):355--363, 1999.

\bibitem{catoni2007pac}
O.\ Catoni,
\textit{PAC-Bayesian Supervised Classification}.
IMS Lecture Notes, 2007.

\bibitem{liu2023pactuning}
Y.\ Liu et al.,
``PAC-tuning: Fine-tuning Pretrained Language Models with PAC-Bayes,''
\textit{arXiv:2310.17588}, 2023.

\bibitem{song2024sparse}
W.\ Song, Z.\ Li, L.\ Zhang, H.\ Zhao, and B.\ Du,
``Sparse is Enough in Fine-tuning Pre-trained Large Language Models,''
\textit{Proceedings of the 41st International Conference on Machine Learning
  (ICML 2024 Spotlight)}, 2024.
arXiv:2312.11875.

\bibitem{takatsu2022vmf}
A.\ Takatsu,
``The $f$-divergence of a von Mises-Fisher Distribution from Another,''
\textit{arXiv:2202.05192}, 2022.

\bibitem{zhang2023adalora}
Q.\ Zhang, M.\ Chen, A.\ Bukharin, N.\ Karampatziakis, P.\ He,
Y.\ Cheng, W.\ Chen, and T.\ Zhao,
``AdaLoRA: Adaptive Budget Allocation for Parameter-Efficient Fine-Tuning,''
\textit{11th International Conference on Learning Representations (ICLR)},
2023. arXiv:2303.10512.

\bibitem{foret2021sam}
P.\ Foret, A.\ Kleiner, H.\ Mobahi, and B.\ Neyshabur,
``Sharpness-Aware Minimization for Efficiently Improving Generalization,''
\textit{9th International Conference on Learning Representations (ICLR)},
2021. arXiv:2010.01412.

\bibitem{bahri2022sharpness}
D.\ Bahri, H.\ Mobahi, and Y.\ Tay,
``Sharpness-Aware Minimization Improves Language Model Generalization,''
\textit{Proceedings of the 60th Annual Meeting of the Association for
  Computational Linguistics (ACL)}, pages 7360--7371, 2022.
DOI:~10.18653/v1/2022.acl-long.508.

\bibitem{vu2020exploring}
T.\ Vu, T.\ Wang, T.\ Munkhdalai, A.\ Sordoni, A.\ Trischler,
A.\ Mattarella-Micke, S.\ Maji, and M.\ Iyyer,
``Exploring and Predicting Transferability across NLP Tasks,''
\textit{Proceedings of the 2020 Conference on Empirical Methods in Natural
  Language Processing (EMNLP)}, 2020.
\url{https://aclanthology.org/2020.emnlp-main.635}.

\bibitem{yang2024unveiling}
H.\ Yang, Y.\ Zhang, J.\ Xu, H.\ Lu, P.\,A.\ Heng, and W.\ Lam,
``Unveiling the Generalization Power of Fine-Tuned Large Language Models,''
\textit{Proceedings of the 2024 Conference of the North American Chapter of the
  Association for Computational Linguistics (NAACL)}, 2024.
arXiv:2403.09162.

\bibitem{tanwar2025domain}
E.\ Tanwar, D.\ Nathani, W.\,Y.\ Wang, and T.\ Chakraborty,
``Understanding the Effects of Domain Finetuning on LLMs,''
\textit{arXiv:2510.09359} [cs.CL], 2025. (Preprint, not peer-reviewed.)

\bibitem{dziugaite2017computing}
G.\ K.\ Dziugaite and D.\ M.\ Roy,
``Computing Nonvacuous Generalization Bounds for Deep (Stochastic) Neural
Networks with Many More Parameters than Training Data,''
\textit{Proceedings of the 33rd Conference on Uncertainty in Artificial
  Intelligence (UAI)}, 2017. arXiv:1703.11008.

\bibitem{liu2019roberta}
Y.\ Liu, M.\ Ott, N.\ Goyal, J.\ Du, M.\ Joshi, D.\ Chen, O.\ Levy,
M.\ Lewis, L.\ Zettlemoyer, and V.\ Stoyanov,
``RoBERTa: A Robustly Optimized BERT Pretraining Approach,''
\textit{arXiv:1907.11692}, 2019.

\bibitem{touvron2023llama}
H.\ Touvron, L.\ Martin, K.\ Stone, P.\ Albert, A.\ Almahairi, Y.\ Babaei,
N.\ Bashlykov, et~al.,
``Llama 2: Open Foundation and Fine-Tuned Chat Models,''
\textit{arXiv:2307.09288}, 2023.

\bibitem{alpaca2023}
R.\ Taori, I.\ Gulrajani, T.\ Zhang, Y.\ Dubois, X.\ Li, C.\ Guestrin,
P.\ Liang, and T.\ B.\ Hashimoto,
``Stanford Alpaca: An Instruction-Following LLaMA Model,''
\textit{GitHub repository}, 2023.
\url{https://github.com/tatsu-lab/stanford_alpaca}.

\end{thebibliography}

\end{document}
